% Agent 02: Book pages 349--352, Section 2.1 (continued)
% Covers equations (2.1.6b) second part through (2.1.22) and Figure 2.1.1
%
% Page 349: eq (2.1.6b) conclusion, "Radiation from a nonrelativistic beam
%           of electrons," eqs (2.1.7)--(2.1.8), Forbidden region (2.1.9a),
%           Photon-discreteness region (2.1.9b)
% Page 350: Large-angle region eqs (2.1.10)--(2.1.13),
%           Small-angle classical region eq (2.1.14)
% Page 351: eqs (2.1.15)--(2.1.19), start of uncertainty-principle region
% Page 352: eqs (2.1.20)--(2.1.22), Figure 2.1.1

%=============================================================================
% [Continuation from Agent 01: large-angle scattering, eq (2.1.6b)]
%=============================================================================

\noindent it swings into a sharp turn around the nucleus; and it emerges
at $t = +\tau/2$ with its initial velocity precisely reversed.
Consequently, $|\Delta v_\omega|$ is equal to~$2v_f$, and $I(\omega)$~is
\begin{equation}
  I(\omega)
  = \frac{32}{3}\,\frac{e^{2}}{c^{3}}
    \left(\frac{Zr_0\,\omega\,c^{2}}{3}\right)^{\!2/3}
  \quad\text{for}\quad
  \frac{v}{b_{s-1}} < \omega < \frac{v_p}{r_p}
  = \frac{c(2Zr_0\,b_{s-1}^{3})^{1/2}}{b^{3}}\,.
  \tag{2.1.6b}
\end{equation}

%=============================================================================
% Radiation from a nonrelativistic beam of electrons
%=============================================================================

\bigskip
\noindent\textit{Radiation from a nonrelativistic beam of electrons}

\medskip

Consider a single ion of charge~$Z$, at rest in the laboratory frame.
Bombard the ion with a monochromatic beam of electrons with speed
$v \ll 1$ and flux (number per unit time per unit area)~$S$.  Examine the
radiation of frequency~$\nu$ (circular frequency $\omega = 2\pi\nu$)
emitted by the electrons as they scatter off the nucleus.
Let $\mathscr{P}_\nu$ be the total power emitted per unit frequency
(ergs~sec$^{-1}$~Hz$^{-1}$); and define the emission cross section per
unit frequency, $d\sigma/d\nu$, by
\begin{equation}
  \mathscr{P}_\nu = (d\sigma/d\nu)\,S\,h\nu.
  \tag{2.1.7}
\end{equation}
The dependence of the emission cross section on photon frequency~$\nu$
and on electron speed~$v$ is most conveniently expressed in terms of
the two dimensionless parameters
\begin{equation}
  \frac{h\nu}{\tfrac{1}{2}m_e v^{2}}\,,
  \qquad
  \frac{\tfrac{1}{2}m_e v^{2}}{2R_y}
  = \left(\frac{v/c}{\alpha Z}\right)^{\!2}.
  \tag{2.1.8}
\end{equation}

On a sheet of paper (Figure~\ref{fig:2.1.1}) plot
$(h\nu)/(\tfrac{1}{2}m_e v^{2})$ vertically, and plot
$(\tfrac{1}{2}m_e v^{2})/(2R_y)$ horizontally.  This 2-dimensional plot
can be split into several different regions, in each of which the details
of the emission process are different.

%--- Forbidden region ---------------------------------------------------
\bigskip
\noindent\textit{Forbidden region}\quad
The individual photons emitted while scattering cannot have energies
greater than the electron kinetic energy---unless the electron gets
captured into a bound state around the ion.  But capture produces
``free-bound'' rather than bremsstrahlung radiation (\S\,2.2).
Consequently, the region
\begin{equation}
  (h\nu)/(\tfrac{1}{2}m_e v^{2}) > 1
  \tag{2.1.9a}
\end{equation}
of Figure~\ref{fig:2.1.1} is a ``forbidden region''; no bremsstrahlung
is emitted in this region; $d\sigma/d\nu$ vanishes in this region.

%--- Photon-discreteness region -----------------------------------------
\bigskip
\noindent\textit{Photon-discreteness region}\quad In the region
\begin{equation}
  \tfrac{1}{3} < (h\nu)/(\tfrac{1}{2}m_e v^{2}) \le 1
  \tag{2.1.9b}
\end{equation}
no more than 3~photons can be emitted by each scattering.  This ``photon
discreteness'' has a significant effect on the emission cross section,
cutting it down toward zero as one approaches the edge of the forbidden
region, $(h\nu)/(\tfrac{1}{2}m_e v^{2}) = 1$.  No classical calculation
can reveal such an effect; the photon-discreteness region must be
analyzed in a quantum mechanical manner; see, e.g., Heitler
\cite{Heitler1954}, and see below.

%=============================================================================
% Large-angle region (page 350)
%=============================================================================

\bigskip
\noindent\textit{Large-angle region}\quad
All small-angle scatterings last too long
($\Delta\tau \sim b/v > b_{s-1}/v$; cf.\ eq.~(2.1.4)) to produce any
radiation of $\tau = 1/\omega < b_{s-1}/v$; cf.\ eq.~(2.1.4).
Consequently, in the region
\begin{equation}
  \frac{h\nu}{\tfrac{1}{2}m_e v^{2}}
  > \frac{h\omega}{\tfrac{1}{2}m_e v^{2}}
  = \left(\frac{\tfrac{1}{2}m_e v^{2}}{Z^{2}R_y}\right)^{\!1/2}
  \tag{2.1.10}
\end{equation}
only large-angle scatterings produce radiation.  This region is shown in
Figure~\ref{fig:2.1.1}.  The power per unit frequency emitted at a fixed
frequency~$\nu$ in this region is
\begin{equation}
  \mathscr{P}_\nu
  = 2\pi \int_0^{b_{\max}} I(\omega)\,S\,2\pi\,b\,db\,,
  \tag{2.1.11a}
\end{equation}
where $I(\omega)$ is given by the large-angle formula~(2.1.6), and
$b_{\max}$ is the ``cutoff''
\begin{equation}
  b_{\max} = (2c^{2}\,Zr_0\,b_{s-1}^{3})^{1/6}
  \tag{2.1.11b}
\end{equation}
in that formula.

Performing the integration, dividing by $S\,h\nu$ to obtain
$d\sigma/d\nu$, and reexpressing the result in terms of fundamental
atomic constants, one obtains
\begin{equation}
  \frac{d\sigma}{d\nu}
  = \frac{2^{1/3}\,16\,\alpha\,c^{2}\,Z^{2}\,r_0^{2}}
         {3^{5/3}\;v^{2}}\,.
\end{equation}

A more exact classical calculation gives a slightly different numerical
coefficient:
\begin{equation}
  \left(\frac{d\sigma}{d\nu}\right)_{\!\text{LA}}
  = \frac{16\pi}{3\sqrt{3}}\;
    \frac{\alpha\,c^{2}\,Z^{2}\,r_0^{2}}{v^{2}}\,.
  \tag{2.1.12}
\end{equation}

\noindent(Here ``LA'' stands for ``large-angle region''.)

In other regions of Figure~\ref{fig:2.1.1}, one gets other formulas for
$d\sigma/d\nu$ (see below).  It is conventional in all regions to lump
the deviations from this large-angle cross section into a correction
$G(\nu,v)$ which is called the ``Gaunt factor'':
\begin{equation}
  G(\nu,v)
  \equiv \frac{(d\sigma/d\nu)}{(d\sigma/d\nu)_{\text{LA}}}\,.
  \tag{2.1.13}
\end{equation}

Thus, in the large-angle region $G = 1$; and in the forbidden region
$G = 0$.  Throughout the large-angle and photon-discreteness portion of
the large-angle region, $G$~remains approximately~1; see
Figure~\ref{fig:2.1.1}.

%=============================================================================
% Small-angle, classical region (pages 350--351)
%=============================================================================

\bigskip
\noindent\textit{Small-angle, classical region}\quad
In the region $\tau = 1/\omega \gg b_{s-1}/v$---i.e.
\begin{equation}
  \frac{h\nu}{\tfrac{1}{2}m_e v^{2}}
  \ll \left(\frac{\tfrac{1}{2}m_e v^{2}}{2R_y}\right)^{\!1/2},
  \tag{2.1.14}
\end{equation}
radiation is produced by small-angle scatterings as well as by
large-angle scatterings.  As a function of impact parameter~$b$, the
energy radiated by a single electron is constant in the large-angle
region, $b < b_{s-1}$ [eq.~(2.1.6b)], and decreases as~$1/b^{2}$ in
the small-angle region, $b > b_{s-1}$ [eq.~(2.1.4)].  Hence, when one
integrates over $2\pi\,b\,db$ to get the power radiated,
$\mathscr{P}_\nu$, one obtains

\smallskip
\noindent(contribution from small angles)
$\displaystyle{}= \ln\!\left(\frac{b_{\max}}{b_{s-1}}\right)$,

\smallskip
\noindent(contribution from large angles) ${}= \text{const}$,

\smallskip
\noindent where $b_{\max}$ is the cutoff in the small-angle region
(eq.~2.1.4):
\begin{equation}
  b_{\max} = v/\omega.
  \tag{2.1.15}
\end{equation}

Since the small-angle contribution is logarithmically dominant, one can
ignore the contribution from large-angle scatterings and calculate
\begin{equation}
  \frac{d\sigma}{d\nu}
  = \frac{\mathscr{P}_\nu}{S\,h\nu}
  = \frac{16\,\alpha\,c^{2}\,Z^{2}\,r_0^{2}}{3\,v^{2}}\,
    \ln\!\left(\frac{b_{\max}}{b_{s-1}}\right).
  \tag{2.1.16}
\end{equation}

\noindent A more exact classical calculation gives a slightly different
argument in the logarithm:
\begin{equation}
  \frac{d\sigma}{d\nu}
  = \frac{16\,\alpha\,c^{2}\,Z^{2}\,r_0^{2}}{3\,v^{2}}\,
    \ln\!\left(\frac{2\,b_{\max}}{\xi\;b_{s-1}}\right),
  \tag{2.1.18}
\end{equation}
where $\xi = 1.781\ldots$ is the ``Euler constant''.  Consequently, the
Gaunt factor for the small-angle, classical region is
\begin{equation}
  G(\nu,v)
  = \frac{\sqrt{3}}{\pi}\,
    \ln\!\left[
      \frac{2}{\xi}
      \left(\frac{\tfrac{1}{2}m_e v^{2}}{h\nu}\right)
      \left(\frac{\tfrac{1}{2}m_e v^{2}}{2R_y}\right)^{\!1/2}
    \right].
  \tag{2.1.19}
\end{equation}

Actually, this small-angle classical result is not valid throughout the
region~(2.1.14).  The uncertainty principle requires a modification at
electron energies $\tfrac{1}{2}m_e v^{2}$ larger than~$Z^{2}R_y$.

%=============================================================================
% Small-angle, uncertainty-principle region (pages 351--352)
%=============================================================================

\bigskip
\noindent\textit{Small-angle, uncertainty-principle region}\quad
Consider an electron with impact parameter~$b$ and speed~$v$.  The
electron is actually not a classical object; rather, it is a
quantum-mechanical wave packet.  To scatter with impact parameter~$b$,
the wave packet (i)~must have transverse dimensions, $\Delta x$, less
than~$b$
\begin{equation}
  \Delta x < b,
\end{equation}
and (ii)~must not spread transversely by more than~$b$ during the
classical scattering time $\Delta t = b/v$
\begin{equation}
  b > \left(\frac{\text{spreading}}{\text{during }\Delta t}\right)
  = \left(\frac{\Delta p_x}{m_e}\right)\Delta t
  > \left(\frac{\hbar}{m_e\,\Delta x}\right)\Delta t
  > \frac{\hbar}{m_e\,v}
  = \frac{b}{m_e\,v}\,\frac{\hbar}{b}\,.
\end{equation}

\noindent Thus, the uncertainty principle (``$\ge$'' in the above
equation) prevents the existence of scatterings with~$b$ less than the
electron de~Broglie wavelength
\begin{equation}
  b_{dB} = \hbar / m_e v.
  \tag{2.1.20}
\end{equation}

If $b_{dB} < b_{s-1}$, this limitation has no effect on small-angle
scatterings; and the small-angle, classical Gaunt factor~(2.1.19) is
valid.  If $b_{dB} > b_{s-1}$, the uncertainty principle comes into
play, and one must use~$b_{dB}$ rather than~$b_{s-1}$ as the lower
limit on the small-angle integral~(2.1.17) for $d\sigma/d\nu$.  The
result is
\begin{equation}
  \frac{d\sigma}{d\nu}
  = \frac{16\,\alpha\,c^{2}\,Z^{2}\,r_0^{2}}{3\,v^{2}}\,
    \ln\!\left(\frac{2\,b_{\max}}{\xi\;b_{dB}}\right),
  \tag{2.1.21}
\end{equation}
corresponding to the Gaunt factor
\begin{equation}
  G(\nu,v)
  = \frac{\sqrt{3}}{\pi}\,
    \ln\!\left(\frac{2\,b_{\max}}{\xi\;b_{dB}}\right)
  = \frac{\sqrt{3}}{\pi}\,
    \ln\!\left[
      \frac{2}{\xi}
      \left(\frac{\tfrac{1}{2}m_e v^{2}}{h\nu}\right)
    \right].
  \tag{2.1.22}
\end{equation}

The ``small-angle, uncertainty-principle region,'' in which this
expression holds, is separated from the ``small-angle, classical
region'' by the line
\begin{equation}
  b_{dB}/b_{s-1}
  = (\tfrac{1}{2}m_e v^{2})/(Z^{2}R_y) = 1.
\end{equation}

\noindent See Figure~\ref{fig:2.1.1}.

From Figure~\ref{fig:2.1.1} it is clear that the uncertainty principle
cannot have any effect on the large-angle region.

%=============================================================================
% Figure 2.1.1
%=============================================================================

\begin{figure}[htbp]
\centering
% \includegraphics[width=\columnwidth]{fig_2_1_1}
\fbox{\parbox{0.9\columnwidth}{\centering
[Figure 2.1.1 placeholder]\\[6pt]
Horizontal axis:
$\bigl(\tfrac{1}{2}m_e v^{2}\bigr)\big/\bigl(Z^{2}R_y\bigr)$.\quad
Vertical axis:
$h\nu\big/\bigl(\tfrac{1}{2}m_e v^{2}\bigr)$.\\[4pt]
Five regions shown:\\[2pt]
(1) Forbidden region ($G = 0$, $d\sigma/d\nu = 0$)
    --- above the line $h\nu/(\tfrac{1}{2}m_ev^2) = 1$.\\
(2) Photon-discreteness region ($G \lesssim 1$)
    --- between $\tfrac{1}{3}$ and~$1$.\\
(3) Large-angle region ($G = 1$)
    --- below discreteness, above
    $h\nu/(\tfrac{1}{2}m_ev^2) =
     [(\tfrac{1}{2}m_ev^2)/(2R_y)]^{1/2}$.\\
(4) Small-angle, classical region:
    $G = \dfrac{\sqrt{3}}{\pi}
    \ln\!\Bigl[\dfrac{4}{\xi}
    \bigl(\dfrac{\frac{1}{2}m_ev^{2}}{h\nu}\bigr)
    \bigl(\dfrac{\frac{1}{2}m_ev^{2}}{2R_y}\bigr)^{\!1/2}\Bigr]$
    --- lower-left.\\[4pt]
(5) Small-angle, U.P.\ (uncertainty-principle) region:
    $G = \dfrac{\sqrt{3}}{\pi}
    \ln\!\Bigl[\dfrac{4}{\xi}
    \bigl(\dfrac{\frac{1}{2}m_ev^{2}}{h\nu}\bigr)\Bigr]$
    --- lower-right.
}}
\caption{Regions and Gaunt factors for bremsstrahlung of frequency~$\nu$
emitted by an electron of kinetic energy~$\tfrac{1}{2}m_e v^{2}$
when it impinges on an ion of charge~$Ze$.}
\label{fig:2.1.1}
\end{figure}
